\documentclass{abntex2}
\usepackage[utf8]{inputenc}
\usepackage[T1]{fontenc}
\usepackage{geometry}
\usepackage{graphicx}
\usepackage{float}
\usepackage{hyperref}
\usepackage[english,indonesian]{babel}

\title{Optimasi Pemrosesan Template LaTeX Menggunakan Pendekatan Berbasis AI}
\author{Academic Support Agent}
\date{\today}

\begin{document}
\maketitle

\begin{abstract}

\end{abstract}

\chapter{PENDAHULUAN}


\chapter{TINJAUAN PUSTAKA}


% Original content (commented out):
% % ------------------------------------------------------------------------
% ------------------------------------------------------------------------
% Modelo de relatório de estágio DAS - UFSC.
% Contribua em https://github.com/gabrielpra1/relatorio-estagio-das
% ------------------------------------------------------------------------
% ------------------------------------------------------------------------

\documentclass[
	% -- opções da classe memoir --
	12pt,				% tamanho da fonte
	openright,			% capítulos começam em pág ímpar (insere página vazia caso preciso)
	oneside,			% Para impressão simples. Para impressão frente e verso, use twoside
	a4paper,			% tamanho do papel. 
	sumario=tradicional,
	% -- opções do pacote babel --
	english,			% idioma adicional para hifenização
	brazil				% o último idioma é o principal do documento
	]{abntex2}


% ---
% Configurações do Preambulo
% ---
\include{preambulo/config}
\include{preambulo/tabela}

\usepackage{blindtext}
\usepackage{lipsum}

% ----
% Início do documento
% ----
\begin{document}

% Retira espaço extra obsoleto entre as frases.
\frenchspacing 

% ----------------------------------------------------------
% ELEMENTOS PRÉ-TEXTUAIS
% ----------------------------------------------------------
% \pretextual

% ---
% Capa
% ---
\include{extra/capa}
% ---

% ---
% Folha de rosto
% (o * indica que haverá a ficha bibliográfica)
% ---
\include{extra/folha_de_rosto}
% ---

% ---
% Inserir folha de aprovação
% ---

% Isto é um exemplo de Folha de aprovação, elemento obrigatório da NBR
% 14724/2011 (seção 4.2.1.3). Você pode utilizar este modelo até a aprovação
% do trabalho. Após isso, substitua todo o conteúdo deste arquivo por uma
% imagem da página assinada pela banca com o comando abaixo:
%
% \includepdf{folhadeaprovacao_final.pdf}
%\include{extra/folha_de_aprovacao}

% ---
% Dedicatória
% ---
%\include{extra/dedicatoria}
% ---

% ---
% Agradecimentos
% ---
\include{extra/agradecimentos}
% ---

%----------------------------------------------------------------%  
% Epígrafe
%----------------------------------------------------------------%  
\include{extra/epigrafe}
%----------------------------------------------------------------%  

%----------------------------------------------------------------%  
% RESUMOS
%----------------------------------------------------------------%  
%\include{extra/resumos}
%----------------------------------------------------------------%  

% Inserir listas. 
% Todas essas listas são opcionais, então caso não queira usar nenhuma, basta comentar a próxima linha.
%\input{listas.tex}

%----------------------------------------------------------------%  
% inserir o sumario
%----------------------------------------------------------------%  
\pdfbookmark[0]{\contentsname}{toc}
\tableofcontents*
\cleardoublepage
% ---



% ----------------------------------------------------------
% ELEMENTOS TEXTUAIS
% ----------------------------------------------------------
\textual

\pagestyle{simple} % Coloca somente o número da página no cabeçalho (Ver seção 7.2 do Manual do Memoir)
\aliaspagestyle{chapter}{simple}% Coloca somente o número da página no cabeçalho das páginas subsequentes do capítulo

% ----------------------------------------------------------
% Introdução (exemplo de capítulo sem numeração, mas presente no Sumário)
% ----------------------------------------------------------
%\chapter*[Introdução]{Introdução}
%\addcontentsline{toc}{chapter}{Introdução}
% ----------------------------------------------------------

%\input{introducao.tex}


% Desenvolvimento
\input{desenvolvimento/capitulo_01.tex}
\input{desenvolvimento/capitulo_02.tex}
\input{desenvolvimento/capitulo_03.tex}
\input{desenvolvimento/capitulo_04.tex}
% ----------------------------------------------------------
% Finaliza a parte no bookmark do PDF
% para que se inicie o bookmark na raiz
% e adiciona espaço de parte no Sumário
% ----------------------------------------------------------
\phantompart

% ----------------------------------------------------------
% ELEMENTOS PÓS-TEXTUAIS
% ----------------------------------------------------------
\postextual
% ----------------------------------------------------------

% ----------------------------------------------------------
% Referências bibliográficas
% ----------------------------------------------------------
\bibliography{bibliografia}
% ----------------------------------------------------------
% Apêndices
% ----------------------------------------------------------

% ---
% Inicia os apêndices
% ---
\begin{apendicesenv}

%Textos elaborados pelo próprio autor, a fim de complementar a sua argumentação

\end{apendicesenv}
% ---


% ----------------------------------------------------------
% Anexos
% ----------------------------------------------------------

% ---
% Inicia os anexos
% ---
%\begin{anexosenv}

% ---
%\chapter{Título Opcional.}
% ---
%Documentos não elaborados pelo autor, utilizados de fundamentação (mapas, leis, estatutos).

%\end{anexosenv}

%---------------------------------------------------------------------
% INDICE REMISSIVO
%---------------------------------------------------------------------
\phantompart
\printindex
%---------------------------------------------------------------------

\end{document}

\end{document}